% =====================================================================
%  preamble.tex  --  packages, layout and macros for the whole report
% =====================================================================
\usepackage[T1]{fontenc}
\usepackage[utf8]{inputenc}
\usepackage[sc]{mathpazo}          % Palatino text + matching maths
\usepackage[scaled=0.88]{helvet}   % scalable sans-serif for labels
\usepackage{courier}
\linespread{1.08}
\usepackage[a4paper,left=2.8cm,right=2.5cm,top=2.7cm,bottom=2.8cm,headheight=15pt]{geometry}
\usepackage{microtype}
\usepackage{amsmath,amssymb,mathtools,bm}
\usepackage{graphicx}
\graphicspath{{figures/}}
\usepackage[dvipsnames,table]{xcolor}
\usepackage{booktabs,array,tabularx,multirow}
\usepackage[font=small,labelfont={bf,sf},labelsep=period,margin=0.6cm,skip=8pt]{caption}
\usepackage{subcaption}
\usepackage{siunitx}
\sisetup{detect-all,range-phrase={--},range-units=single,separate-uncertainty=true,per-mode=symbol}
\DeclareSIUnit\ohmpersec{\ohm\per\second}
\usepackage{enumitem}
\setlist{itemsep=0.25em,topsep=0.5em}
\usepackage{needspace}
\usepackage{float}
\usepackage{placeins}
\usepackage{etoolbox}
\renewcommand{\topfraction}{0.9}
\renewcommand{\bottomfraction}{0.8}
\renewcommand{\textfraction}{0.07}
\renewcommand{\floatpagefraction}{0.8}
\setcounter{topnumber}{3}
\setcounter{bottomnumber}{2}
\setcounter{totalnumber}{4}
\usepackage{tikz}
\usetikzlibrary{arrows.meta,positioning,calc,shapes.geometric,shapes.misc,fit,backgrounds,decorations.pathreplacing}
\usepackage[siunitx,RPvoltages]{circuitikz}
\usepackage[most]{tcolorbox}
\usepackage{titlesec}
\usepackage{fancyhdr}
\usepackage[lof,lot]{etoc}
\usepackage[numbers,sort&compress]{natbib}
\usepackage[hidelinks]{hyperref}
\usepackage[nameinlink,capitalise,noabbrev]{cleveref}

% ---------------- colours ---------------------------------------------
\definecolor{accent}{HTML}{7A2E22}     % deep maroon: chapter numbers, rules
\definecolor{ink}{HTML}{1F2A44}        % links
\definecolor{soft}{HTML}{5B5B5B}
\definecolor{boxbg}{HTML}{F3F3F3}
\definecolor{keybg}{HTML}{EDEDED}
\definecolor{boxrule}{HTML}{9A9A9A}
\definecolor{boxhead}{HTML}{5A5A5A}
\definecolor{upcol}{HTML}{1F6F3F}      % switch-on / recovers
\definecolor{dncol}{HTML}{B2182B}      % switch-off / tipped
\hypersetup{colorlinks=true,linkcolor=ink,citecolor=ink,urlcolor=ink,
  pdftitle={Rate-Induced Tipping in a Wien-Bridge Oscillator},
  pdfauthor={Dhruv Kumar Solanki}}

% ---------------- project-wide facts -----------------------------------
\newcommand{\reporttitle}{Rate-Induced Tipping in a Wien-Bridge Oscillator}
\newcommand{\repourl}{https://github.com/dhruv-solanki-2564/MSP-1_R_TIPPING_WEIN_BRIDGE}

% ---------------- headings ---------------------------------------------
\titleformat{\chapter}[display]
  {\normalfont\bfseries}
  {\color{accent}\Large\scshape\chaptertitlename\ \thechapter}
  {0.4em}
  {\Huge\raggedright}
  [\vspace{0.6em}{\color{accent}\titlerule[1.2pt]}]
\titlespacing*{\chapter}{0pt}{-1.2em}{1.6em}
\titleformat{\section}{\normalfont\Large\bfseries}{\color{accent}\thesection}{0.8em}{}
\titleformat{\subsection}{\normalfont\large\bfseries}{\color{accent}\thesubsection}{0.8em}{}
\titleformat{\paragraph}[runin]{\normalfont\bfseries}{}{0pt}{}[.]
\titlespacing*{\section}{0pt}{1.6em}{0.7em}
\titlespacing*{\subsection}{0pt}{1.2em}{0.5em}
\setlength{\parindent}{0pt}
\setlength{\parskip}{0.65em plus 0.15em minus 0.1em}

% ---------------- running header ---------------------------------------
\pagestyle{fancy}
\fancyhf{}
\renewcommand{\chaptermark}[1]{\markboth{#1}{}}
\fancyhead[L]{\small\itshape\color{soft}\nouppercase{\leftmark}}
\fancyhead[R]{\small\color{soft}M.Sc.\ Project Stage~I}
\fancyfoot[C]{\small\thepage}
\renewcommand{\headrulewidth}{0.4pt}
\fancypagestyle{plain}{\fancyhf{}\fancyfoot[C]{\small\thepage}\renewcommand{\headrulewidth}{0pt}}

% ---------------- per-chapter lists of figures and tables ---------------
% \chapterlists prints, right under a chapter title, the figures and tables
% of that chapter with clickable page references.
\makeatletter
\newcommand{\chapterlistline}{\noindent\makebox[2.6em][l]{\etocnumber}%
  \parbox[t]{\dimexpr\linewidth-5.2em}{\raggedright\etocname\strut}\hfill\etocpage\par\vspace{1.5pt}}
\renewcommand{\etoclocallistoffigureshook}{\etocsetstyle{section}{}{}{\chapterlistline}{}}
\renewcommand{\etoclocallistoftableshook}{\etocsetstyle{section}{}{}{\chapterlistline}{}}
% optional argument: f = figures only, ft = figures and tables (default), t = tables only
\newcommand{\chapterlists}[1][ft]{%
  \begingroup\small
  \etocsettocstyle{}{}%
  \begin{tcolorbox}[enhanced,colback=white,colframe=boxrule,boxrule=0.4pt,arc=1.5pt,
      left=8pt,right=8pt,top=5pt,bottom=4pt,before skip=0pt,after skip=1.4em]
  \def\chl@arg{#1}\def\chl@f{f}\def\chl@t{t}%
  \ifx\chl@arg\chl@t\else
    {\sffamily\bfseries\footnotesize\color{boxhead}FIGURES IN THIS CHAPTER}\par\vspace{3pt}
    \locallistoffigures
  \fi
  \ifx\chl@arg\chl@f\else
    \ifx\chl@arg\chl@t\else\vspace{5pt}\fi
    {\sffamily\bfseries\footnotesize\color{boxhead}TABLES IN THIS CHAPTER}\par\vspace{3pt}
    \locallistoftables
  \fi
  \end{tcolorbox}
  \endgroup}
\makeatother

% ---------------- boxes -------------------------------------------------
% "Why this step?" side note used inside derivations
\newtcolorbox{whybox}[1][Why this step?]{enhanced,breakable,colback=boxbg,colframe=boxrule,
  boxrule=0pt,leftrule=2.5pt,arc=0pt,left=9pt,right=9pt,top=6pt,bottom=6pt,
  fontupper=\small,before skip=0.9em,after skip=0.9em,
  title={#1},coltitle=boxhead,fonttitle=\sffamily\bfseries\small,
  attach title to upper={\ \ },titlerule=0pt,colbacktitle=boxbg}
% key result / take-away box
\newtcolorbox{keybox}[1][]{enhanced,breakable,colback=keybg,colframe=boxrule,boxrule=0.5pt,arc=2pt,
  left=10pt,right=10pt,top=7pt,bottom=7pt,before skip=1em,after skip=1em,#1}
% end-of-chapter summary
\newtcolorbox{learned}[1][What this stage taught us]{enhanced,breakable,colback=boxbg,colframe=boxrule,
  boxrule=0.6pt,arc=2pt,left=10pt,right=10pt,top=7pt,bottom=7pt,before skip=1.2em,after skip=1em,
  title={#1},fonttitle=\sffamily\bfseries\small,colbacktitle=boxhead,coltitle=white}

% ---------------- maths shorthands --------------------------------------
\newcommand{\dd}{\mathrm{d}}
\newcommand{\RF}{R_{F}}
\newcommand{\Ri}{R_{i}}
\newcommand{\Gc}{G_{c}}
\newcommand{\wo}{\omega_{0}}
\newcommand{\Vout}{V_{\mathrm{out}}}
\newcommand{\Asat}{A_{\mathrm{sat}}}
\newcommand{\Ron}{R_{\mathrm{on}}}
\newcommand{\Roff}{R_{\mathrm{off}}}
\newcommand{\Lon}{\Lambda_{\mathrm{on}}}
\newcommand{\Loff}{\Lambda_{\mathrm{off}}}
\newcommand{\Gp}{\Gamma_{+}}
\newcommand{\Gm}{\Gamma_{-}}
\newcommand{\lamc}{\lambda_{c}}
\newcommand{\kOhm}{\si{\kilo\ohm}}
\crefname{equation}{Eq.}{Eqs.}
\Crefname{equation}{Equation}{Equations}
\crefname{figure}{Fig.}{Figs.}
\Crefname{figure}{Figure}{Figures}
\crefname{chapter}{Chapter}{Chapters}
\crefname{section}{Section}{Sections}
\crefname{table}{Table}{Tables}
